\chapter{Experimental Methodology}\label{chap:methodologie}

\section{Introduction}
\TODO{Chapter introduction}

\section{Task Selection and Feature Availability}
% Retrospective log classification. Table: features available before the decision / after the decision.
% Justify whether NAT-related features and traffic volume features (Bytes, Packets, Elapsed Time) are used.
\TODO{Selected task and feature availability table}

\section{Preprocessing}
% Duplicate handling (decision and justification), log1p transformation, standardization (KNN, SVM),
% port encoding; everything must be included in a Pipeline fitted on the training data only (no data leakage).
\TODO{Preprocessing pipeline with a diagram (draw.io) and justification for each step}

\section{Data Splitting}
% Grouped splitting: identical feature vectors must remain on the same side. Class counts (training / test),
% random seed, number of rare-class examples in each split.
\TODO{Principle, verified counts, and data splitting diagram}

\section{Experimental Protocol}
% Compared models, hyperparameter grids (table), cross-validation (K >= 5), systematic search,
% evaluation metrics, final test set kept untouched until model selection is complete.
\TODO{Complete experimental protocol and hyperparameter grid table}

\section{Conclusion}
\TODO{Chapter conclusion and transition to Chapter~\ref{chap:resultats}}